\scanpage{51}
Caractérisation de $W_\infty$ (\unsure{Hot}-équivalences) dans $\mathrm{Cat}$

$\overline{\underline{\mathrm{XII}}}$

\unsure{Axiomes} $W(4\ldots)$ p. 13-14

$\overline{\underline{\mathrm{I}}}$ ``Caractérisation de $W_\infty$'' p. 1

$\overline{\underline{\mathrm{II}}}$ Perplexités sur les axiomes et les notions (de $W$-fibration notamment) p. 18-22

$\overline{\underline{\mathrm{III}}}$ Morphismes $W$-propres, $W$-lisses, $W$-fibrants p. 22-52

$\overline{\underline{\mathrm{IV}}}$ Digression sur l'axiome $W(4)$ p. 53-58

$\overline{\underline{\mathrm{V}}}$ Sommes amalgamées et carrés \inserted{\unsure{homotopiquement}} \unsure{cocartésiens} dans $(\mathrm{Cat})$ \inserted{(cf notamment th. 4 (p. 145) et th. \unsure{6|5} (p. \unsure{157|177}))} p. 59-169

$\overline{\underline{\mathrm{VI}}}$ \inserted{\S\S\,1,2} \deleted{Les \unsure{digressions|blessures} \unsure{ici}} Critère de $W_\infty$-asphéricité \deleted{\unsure{en}} en termes de la $0$-connexité des $\Omega^{\unsure{i}}(X,\unsure{x|n})$. \unsure{Suite}. p. 171-\deleted{\unsure{191}}190

\deleted{Digression sur la philosophie des \unsure{fibres} homotopiques etc \textemdash\ p. 191-26\unsure{4|9}}

\deleted{Cofibres $W$-homotopiques dans $\mathrm{Cat}$ : p. 227-239}

\deleted{Retour sur les axiomes $W(7)$, $W(7\mathrm{bis})$ \enspace 265-275}

\unsure{Suite|Sortie} des carrés $W$-cartésiens et cocartésiens \enspace 258-264

\S\,3 Sur la philosophie des fibres homotopiques (\unsure{divagations} initiales) \enspace 191-207

\S\,4 Fibres $W$-homotopiques (1) (notamment a) b) c) p. 215, 218) \enspace 207-218

\S\,5 $W$-types d'homotopie relatifs sur localiseur \unsure{constant} \enspace 218-227

\S\,6 Cofibres $W$-homotopiques dans $\mathrm{Cat}$ \enspace 227-240

\S\,7 Fibres $W$-homotopiques (2) ($+$lemme du \unsure{foncteur|fondeur} adjoint) (Théorie des carrés $W$-cartésiens, p. 258-264) \enspace 240-264

\S\,8 Retour \unsure{sur} les axiomes $W(7)$, $W(7\mathrm{bis})$ \enspace 265-275

\uline{Appendices :} $\overline{\underline{\mathrm{I}}}$) Contre-exemples \deleted{\unsure{réfutés|réfutatifs}} : la notion de flèches tot. $W$-asph : \hfill A1-A16

\diagram{\begin{tikzcd}[ampersand replacement=\&,column sep=small,row sep=tiny]
\& \& \bullet\\
\Psi={} \& \bullet \arrow[ur,dash] \arrow[dr,dash] \& \& \text{n'est \uline{pas} \unsure{tot.|LA} $W_2$-asph.}\\
\& \& \bullet
\end{tikzcd}}

$\overline{\underline{\mathrm{II}}}$) Recollement d'un ouvert et d'un fermé comme \inserted{\unsure{$W_\infty$-|théo-}} \unsure{somme} \unsure{amalgamée} (conjectural) \hfill A17-A18

% archive sheet number: 51
